%!TeX root=../tese.tex
%("dica" para o editor de texto: este arquivo é parte de um documento maior)
% para saber mais: https://tex.stackexchange.com/q/78101

%% ------------------------------------------------------------------------- %%

% "\chapter" cria um capítulo com número e o coloca no sumário; "\chapter*"
% cria um capítulo sem número e não o coloca no sumário. A introdução não
% deve ser numerada, mas deve aparecer no sumário. Por conta disso, este
% modelo define o comando "\chapter**".
\chapter**{Introdução}
\label{cap:introducao}

\enlargethispage{.5\baselineskip}


A Região Metropolitana de São Paulo possui cerca de 20,5 milhões de habitantes, sendo 11 milhões apenas da cidade de São Paulo. Toda essa população enfrenta, diariamente, o estresse do transporte, seja ele coletivo ou individual, gastando, em média, 2h47 por dia ao optar pelo transporte coletivo, em sua maioria, ônibus \citep{redenossa2024}. 

A SPTrans, empresa responsável pelo transporte municipal de São Paulo, disponibiliza na web dados sobre as linhas e quantidade de passageiros transportados. De acordo com o Relatório Integrado de Administração da Empresa de 2024, foram transportados, em média, 7,13 milhões de passageiros por dia útil, através de 1.320 linhas, cobrindo 4,7 mil km, com uma frota de 13.277 veículos \citep{sptrans2024}. O sistema de transporte que a empresa fornece, porém, está longe de atender à crescente demanda da população.

Atento à problemática do transporte paulista, \citet{tcc:Andre} propõe uma análise da qualidade e confiabilidade dos dados GTFS (General Transit Feed Specification) disponibilizados pela SPTrans. Utilizando ferramentas como Node.js e MongoDB, seu trabalho identifica falhas, inconsistências e características operacionais peculiares no sistema de ônibus da cidade, como variações de frequência, duplicidade de paradas e velocidades inconsistentes. A partir disso, valida o uso dos dados para análises críticas sobre a operação e o planejamento da rede, além de propor métricas para avaliação da qualidade do serviço, concluindo que, apesar de possuir algumas inconsistências, os dados disponibilizados pela SPTrans podem ser usados para estudos aprofundados do transporte público operado pela empresa na cidade, seguindo a recomendação de estruturação de dados do GTFS. No entanto, o estudo mantém o foco em análises baseadas em registros operacionais e não se debruça sobre a estrutura territorial do sistema ou sobre sua cobertura em relação à população.

Por outro lado, \citet{tese:sandro} avalia a robustez da rede de transporte público da cidade de São Paulo por meio da modelagem da rede como um grafo complexo, utilizando dados do GTFS, focando em compreender a estrutura fundamental de conectividade da rede e a resiliência do sistema de transporte público de São Paulo por meio de ferramentas da ciência de redes complexas. Seu trabalho se destaca pela aplicação de simulações de ataques direcionados e aleatórios à rede, com o objetivo de compreender sua vulnerabilidade estrutural. Foram analisadas métricas como o número de links, o grau máximo dos nós, a quantidade de componentes e o comprimento médio dos caminhos, revelando pontos críticos de conexão, padrões de conectividade e a resiliência da rede sob diferentes cenários, concluindo que a rede é robusta contra falhas aleatórias, e vulnerável à remoção de pontos por onde passam muitas linhas, como terminais de ônibus. Embora a abordagem ofereça uma leitura da estrutura e conectividade da rede, sua análise é concentrada em aspectos estruturais e não considera a dimensão espacial detalhada da infraestrutura de paradas ou a articulação com a malha urbana. Sandro demonstra que os dados públicos disponibilizados pela SPTrans possibilitam análises profundas e sofisticadas da rede de transporte coletivo do município de São Paulo. Esse exemplo reforça a viabilidade de empregar essas bases de dados, amplamente acessíveis à população, para investigações de alto nível que contribuam com o planejamento urbano e a formulação de políticas públicas que beneficiem a população e o município.

De forma complementar, \citet{tcc:Danilo} realiza a modelagem do transporte público de ônibus por meio de grafos não orientados, mapeando as principais rotas e suas conexões entre regiões político-administrativas da cidade (zonas). Seu trabalho propõe uma metodologia para representar as rotas em grafos, permitindo análises de caminhos mínimos, fluxos e identificação de regiões com maior concentração de linhas, levando em consideração os dados do transporte em horários de pico da manhã e tarde. Ao atribuir pesos aos nós com base no número de linhas conectadas, Silva infere padrões de demanda, apontando, por exemplo, a saturação nas zonas leste e sul da cidade. O trabalho representa um avanço na análise de demanda a partir da estrutura das linhas, porém, ainda se restringe à malha de ônibus e não aprofunda a integração entre diferentes modos de transporte nem a distribuição geoespacial das paradas, podendo ser melhorado com a adição de dados geográficos complementares.

Dado que esses estudos iniciais estabeleceram bases sólidas para a modelagem e análise da rede de transporte paulistana, este trabalho pretende contribuir para a gestão da mobilidade no município de São Paulo trazendo uma análise que contemple os períodos pré, durante e pós-pandemia, evidenciando mudanças e identificando pontos de melhoria a serem estudados no planejamento estratégico da malha viária da cidade, levando em conta o histórico de dados da SPTrans de 2015 a 2025. Para isso, propomos uma análise territorial detalhada da infraestrutura do transporte público, com foco na distribuição geoespacial das paradas e sua articulação com outras bases urbanas, como estações de metrô e trem, e zonas demográficas. A partir da integração de múltiplas fontes de dados, criamos uma base unificada e atualizada, que permite analisar a eficiência e viabilidade da malha de ônibus municipais ao longo dos anos, por meio de análise de redes, com a criação de grafos, e análises estatísticas descritivas, com o objetivo de elucidar e divulgar a condição e qualidade do transporte público em São Paulo. Assim, amplia-se a compreensão da estrutura do transporte público em São Paulo, oferecendo uma percepção mais atual e espacialmente detalhada da malha de mobilidade urbana.

Os objetivos específicos incluem a criação de um grafo com a rede de ônibus municipais da cidade de São Paulo, a identificação de rotas críticas, com alta quantidade de passageiros, a análise das implicações de novas estações de metrô, a análise estatística descritiva e a comparação entre a oferta de transporte e a demanda populacional. Por meio dessas análises, pretendemos investigar a eficiência e sugerir insights para o estudo de melhorias na estrutura atual, oferecendo dados e análises para o planejamento urbano e o aperfeiçoamento do transporte na cidade. Para isso, buscamos coletar, tratar e integrar dados de diversas fontes, como o GTFS e o histórico de passageiros por linha, fornecido pela SPTrans, e dados populacionais do IBGE, a fim de construir um modelo matemático da malha viária.

\citet{SUGUIY2020938} definem o índice ESR (Efficiency/Satisfaction Ratio) como uma medida que busca equilibrar duas dimensões fundamentais na avaliação de sistemas de transporte urbano: a eficiência operacional e a satisfação dos usuários. Segundo os autores, um sistema é eficiente quando não é possível obter uma saída sem, ao menos, uma entrada, e não mais do que isso, enquanto a satisfação representa o grau em que as necessidades dos usuários são atendidas. Priorizar apenas a eficiência pode comprometer a qualidade do serviço, tornando-o socialmente insustentável e, por outro lado, focar exclusivamente na satisfação pode elevar os custos operacionais e prejudicar a viabilidade econômica, mais especificamente, o lucro. No entanto, segundo \citet{schwartzman1975}, o mercado não pode ser modelo para gestão de bens públicos, apontando que a lógica do lucro leva à ênfase em atividades rentáveis, deixando de lado necessidades sociais relevantes, e que a racionalidade política (controle social, participação, legitimidade) difere da racionalidade econômica do mercado, reforçando que o transporte público deve ser analisado pelos seus impactos sociais e a satisfação dos usuários, e não pela rentabilidade.

Neste trabalho, então, a eficiência no transporte público é entendida como a capacidade em equilibrar demanda e oferta, ou seja, a capacidade do sistema satisfazer adequadamente a demanda da população, oferecendo cobertura territorial, frequência, conforto, tempo de deslocamento e custos compatíveis com as necessidades dos usuários, resultando em maior satisfação e adesão da população ao uso de transporte público, sem que haja gastos governamentais em excesso, isto é, operando as rotas de modo a atender a demanda populacional. 

O tratamento e a análise dos dados foram realizados em um notebook Dell G5, equipado com um processador Intel(R) Core(TM) i5-10500H CPU @ 2.50GHz e 16,0 GB de memória RAM. Essa configuração foi suficiente para suportar as etapas de pré-processamento, modelagem e análise estatística dos dados em ambiente local. O código-fonte desenvolvido para a coleta, tratamento e análise dos dados apresentados neste trabalho está disponível publicamente no GitHub, no seguinte repositório: \url{https://github.com/fran-dea/public-transport-sptrans}. A disponibilização visa garantir a reprodutibilidade dos resultados e facilitar eventuais extensões futuras da pesquisa.